\chapter{Introducción}

[Placeholder]

Aquí se motiva el trabajo, se plantea y define el problema, se deja claro cuales son los objetivos (general del proyecto, si correspondiese o si está inmerso en un proyecto de mayor alcance, y los específicos), se plantean los resultados esperados, se establecen resumidamente las conclusiones y se describe la
organización general del documento.

Este y los siguientes capítulos pueden incluir referencias, algunos ejemplos de referencias pueden ser \cite{CitekeyArticle}, \cite{CitekeyBook}, \cite{CitekeyInproceedings}, \cite{CitekeyManual} y \cite{CitekeyMisc}. Se sugiere citar usando el formato APA.

[Que poner]

Esta introccion para mi amerita incluir aspectos tales como el enfoque principal de investigacion del proyecto tales como es el analisis de textos usando PLN y las dificultades que dispone de por si y las adicionales resultantes de trabajar con mensajes del entorno de redes sociales.
Enfoques por lo que va a ir esta tesis, trabajar con diferentes tipos de prepocesamientos (conbinaciones de los mismos), diferentes modelos de procesamientos de textos para devolver que temas fueron devueltos.

Y establecer milestones sobre que es posible utilizar en el contexto de estudio de mensajes en contextos en redes sociales en la region rioplatense.

\section{Problema y motivación}

Poner algo sobre el flujo de mensajes y como es demasiado para que una persona lo vea, impulsar la necesidad de tener herramientas de estas caracteristicas, 

\section{Objetivos}

Investigar como hacer con el lenguaje ruidoso de redes sociales y como tratarlo, que modelos los pueden procesar mejor estos textos y como puedo efectivamente tratar de procesar el flujo de los mismo y tener una idea aproximada de lo hablado 
